\chapter{Refinement and Scoring}
\label{ch:scoring}

Buckets that survive pruning (Chapter~\ref{ch:pruning}) are handed to \emph{refinement}: compute
an exact similarity score, and the precise matched sub-interval, for the bucket. This chapter
covers all four scoring metrics (\S\ref{sec:metric-axis}), the two-pointer sliding-window sweep
that underlies \code{Containment}/\code{Jaccard} scoring, the LCS-based alternative, and
\code{match\_rest}, the procedure that ties refinement into a best/second-best search over all of
a read's surviving buckets.

\section{The \code{diff\_hist} bookkeeping device}
\label{sec:scoring-sliding-window}

Every sliding-window sweep in this chapter maintains a shared per-read map \code{diff\_hist}:
\code{hash} $\to$ \code{remaining budget}, initialized once per read from each seed's
\code{occs\_in\_p} (how many times that $k$-mer occurs in the read). As the sweep's right (or
left) boundary crosses a reference position whose $k$-mer hash is a read seed, the corresponding
budget entry is decremented (window grows) or incremented (window shrinks) by exactly one, and
the \emph{sign} of the result after a decrement determines whether that occurrence is
countable as part of the intersection:
\begin{itemize}[nosep]
  \item Decrementing (window grows, position enters): if the value \emph{after} decrementing is
    still $\ge 0$, this occurrence is a legitimate, not-yet-used match --- count it, and record
    that it was counted (so the corresponding increment, when the window later shrinks past it,
    knows to un-count it).
  \item Incrementing (window shrinks, position leaves): if the value \emph{after} incrementing is
    $\ge 1$, this position had previously been counted --- un-count it.
\end{itemize}
This lets a sliding window maintain an exact intersection-with-multiplicity count in $O(1)$
amortized per boundary move, correctly handling repeated $k$-mers (a hash occurring 3 times in
the read can be ``used up'' by at most 3 reference positions inside the current window, no
matter how many more copies of that hash the reference bucket actually contains) --- this is the
\emph{multiset} intersection a containment/Jaccard estimate needs (\S\ref{sec:fracminhash-theory}),
not a plain set intersection.

\begin{notebox}
Because a window can shrink and regrow over the same reference position multiple times as the
sweep's left boundary advances, \code{diff\_hist} values legitimately go negative
\emph{transiently} (a decrement can push a shared hash's budget below what's currently
``available'' if multiple reference positions inside the window share that hash but the read only
has fewer occurrences) --- this is expected and exactly what the $\ge 0$ / $\ge 1$ sign checks are
designed to handle correctly; do not add an assertion that rejects negative values.
\end{notebox}

\section{\code{SHMapper}'s own \code{bestFixedLength}: the main refinement sweep}
\label{sec:best-fixed-length}

This is the function that actually refines \code{Containment}/\code{Jaccard} buckets in the main
mapping path (there is a second, differently-shaped function with the same name used only by the
ground-truth diagnostic path, \S\ref{sec:matcher-best-fixed-length} --- keep them distinct).

\paragraph{Input.} The reference segment (its full sketch), a window \code{[from, to)} to sweep
(the bucket's own span, from \code{buckets.begin}/\code{end}), the read's seed lookup table
(hash $\to$ \code{Seed}), \code{diff\_hist} (mutable, shared across the whole read's refinement
calls), the read length minus $k$ (used as the PAF end-coordinate, not a window bound), the window
\emph{width} bound $m$ (confusingly, at the call site this is actually the read's containing
$\ell$/\code{lmax} value, not the sketch size --- see the note below), and the metric
(\code{Containment} or \code{Jaccard}).

\paragraph{Output.} A \code{Mapping} (\S\ref{sec:type-mapping}): the best-scoring sub-window
found anywhere inside \code{[from, to)}.

\begin{algobox}{title={\code{bestFixedLength}(segm, from, to, seed\_table, diff\_hist, p\_sz, m, metric)}}
\begin{verbatim}
t := segm.kmers
l := max(from, 0)
r := l
end := min(len(t), to)
best := Mapping::default()          # score = -1, beaten by any real candidate

while l < end:
    while r < end:
        # window-width check, in whichever coordinate space abs_pos selects
        if abs_pos:  if not (t[r].r < t[l].r + m): break
        else:        if not (r < l + m):           break
        if t[r].h in seed_table:
            same_strand += codirection(seed_table[t[r].h].kmer, t[r])
            if (--diff_hist[t[r].h]) >= 0:
                intersection += 1
        r += 1
    # r is now exactly one past the last position included in this window
    # for the CURRENT l -- this is the point at which it is safe to score,
    # since r has already been advanced past every included element.
    s_kmers := r - l
    score := mappingScore(intersection, m, s_kmers, metric)       # 10.2 below
    if l < r and score > best.score():
        best.update(p_start=0, p_end=p_sz-1, t_l=t[l].r, t_r=t[r-1].r,
                    segm=segm, intersection=intersection, score=score,
                    same_strand=same_strand, span=t[r-1].r - t[l].r + 1)
    if t[l].h in seed_table:
        same_strand -= codirection(seed_table[t[l].h].kmer, t[l])
        if (++diff_hist[t[l].h]) >= 1:
            intersection -= 1
    l += 1
return best
\end{verbatim}
\end{algobox}

\paragraph{\code{mappingScore}:}
\[
  \mathrm{mappingScore}(\mathit{intersection}, m, s\_sz, \mathtt{Jaccard}) =
    \frac{\mathit{intersection}}{m + s\_sz - \mathit{intersection}}, \qquad
  \mathrm{mappingScore}(\mathit{intersection}, m, s\_sz, \mathtt{Containment}) =
    \frac{\mathit{intersection}}{m}.
\]
These are exactly the FracMinHash Jaccard/containment estimators of
\S\ref{sec:fracminhash-theory}, with $m$ the read's sketch size and $s\_sz$ the current window's
$k$-mer count.

\begin{notebox}
\label{sec:m-vs-psz-confusion}
At the actual call site (\S\ref{sec:find-best-mapping}), this function's \code{p\_sz} parameter
receives \emph{read length minus $k$}, and its \code{m} parameter receives $\mathit{lmax}$ (the
bucket half-length $\ell$), \emph{not} the read's sketch size in either case, despite parameter
names that suggest otherwise at a casual read. This is not a bug --- $\ell$ was chosen
(\S\ref{sec:bucket-halflen}) to equal the read's sketch size in the default (non-\code{abs\_pos})
configuration, so the values coincide numerically for the common case, but a from-scratch
implementation should pass the value that is actually semantically the ``window-width bound and
Jaccard/containment denominator'' at this call site, not necessarily whatever local variable
happens to be named \code{m} in a hasty transcription. Name the parameter for what it does,
not what one caller happens to pass.
\end{notebox}

\paragraph{Why scoring happens \emph{after} the inner \code{while} loop exits, not inside it.}
This is the single most important structural detail in this function, and the source of a real
bug in the ground-truth-only sibling function discussed in \S\ref{sec:matcher-best-fixed-length}:
after the inner loop's normal exit, \code{r} is \emph{exactly one past} the last position that
was ever included in the window for the current \code{l} (standard half-open-range convention:
the loop only stops once its own condition fails, having already processed and advanced past
every element that satisfied it). So \code{t[r-1]} unambiguously means ``the last element
actually in the window'', consistent with what \code{intersection} counts. Any refinement sweep
that instead scores \emph{mid-loop} (before the equivalent exit has happened) must use the
\emph{current}, not-yet-advanced index for both the score's span \emph{and} its use of
\code{intersection}, or the two will disagree --- see \S\ref{sec:matcher-best-fixed-length} for a
concrete case where the reference implementation gets exactly this wrong.

\section{\code{findBestMapping}: dispatching on metric}
\label{sec:find-best-mapping}

\begin{algobox}{title={\code{findBestMapping}(buckets, loc, content, seed\_table, diff\_hist, p\_sz, m, lmax, sh, metric, k)}}
\begin{verbatim}
switch metric:
  case bucket_SH:
    # Degenerate/cheap: reuse the seed-heuristic bound itself as the score,
    # with the bucket's own tracked position range as the matched interval.
    # No refinement sweep at all -- O(1) given the bucket's content.
    best := Mapping(p_start=0, p_end=p_sz-1, t_l=content.r_min, t_r=content.r_max,
                    segm=ref_segment(loc.segm_id), intersection=content.matches,
                    score=sh, same_strand=content.codirection,
                    span=content.r_max - content.r_min)

  case bucket_LCS:
    lcs_len := lcs(ref_segment(loc.segm_id).kmers, buckets, loc, seed_table)   # 10.4
    score := lcs_len / m
    best := Mapping(p_start=0, p_end=p_sz-1, t_l=content.r_min, t_r=content.r_max,
                    segm=ref_segment(loc.segm_id), intersection=content.matches,
                    score=score, same_strand=content.codirection,
                    span=content.r_max - content.r_min)

  case Containment or Jaccard:
    best := bestFixedLength(ref_segment(loc.segm_id), buckets.begin(loc), buckets.end(loc),
                             seed_table, diff_hist, p_sz - k, lmax, metric)

best.set_bucket(loc)
best.set_sh(sh)
return best
\end{verbatim}
\end{algobox}

\section{The \code{bucket\_LCS} metric: longest common subsequence chain score}
\label{sec:lcs-implementation}

\paragraph{Motivation.} Containment/Jaccard scoring treats every matched $k$-mer equally,
regardless of whether the matches are mutually \emph{co-linear} (consistent with one real
alignment) or scattered (an artifact of repeats). \code{bucket\_LCS} instead scores a bucket by
how long a co-linear chain of matches it contains, per \S\ref{sec:lis-theory}.

\subsection{\code{lcs\_get\_ppos\_in\_t}: collecting the position sequence}

\begin{algobox}{title={\code{lcs\_get\_ppos\_in\_t}(t, buckets, loc, seed\_table)}}
\begin{verbatim}
ppos_in_t := []
for l in buckets.begin(loc) .. min(len(t), buckets.end(loc)) - 1:
    if t[l].h in seed_table:
        for ppos in seed_table[t[l].h].pmatches:      # Section 4.4: all read positions sharing this hash
            ppos_in_t.append(ppos)
return ppos_in_t
\end{verbatim}
\end{algobox}

This produces, in increasing reference-position order, the sequence of read-side positions that
each reference $k$-mer inside the bucket matches (with repetition, if a hash occurs multiple times
in the read). A long increasing run in this sequence is a long co-linear chain.

\subsection{\code{lcs\_get\_lis}: patience-sorting LIS length}
\label{sec:lcs-lis}

Exactly the classic $O(n\log n)$ algorithm of \S\ref{sec:lis-theory}:

\begin{algobox}{title={\code{lcs\_get\_lis}(ppos\_in\_t)}}
\begin{verbatim}
tails := []                       # tails[i] = smallest tail of an increasing run of length i+1
for x in ppos_in_t:
    lo, hi := -1, len(tails)
    while lo + 1 < hi:
        mid := (lo + hi) / 2
        if tails[mid] < x: lo := mid
        else:              hi := mid
    if hi == len(tails): tails.append(x)
    else:                 tails[hi] := x
return tails                       # len(tails) is the LIS length
\end{verbatim}
\end{algobox}

\subsection{\code{lcs}: combining forward and reverse}

\begin{algobox}{title={\code{lcs}(t, buckets, loc, seed\_table)}}
\begin{verbatim}
ppos_in_t := lcs_get_ppos_in_t(t, buckets, loc, seed_table)
forward_len := len(lcs_get_lis(ppos_in_t))
reverse(ppos_in_t)
reverse_len := len(lcs_get_lis(ppos_in_t))
return max(forward_len, reverse_len)
\end{verbatim}
\end{algobox}

Running LIS on the \emph{reversed} position sequence finds the longest \emph{decreasing} run of
the original --- i.e., the longest chain co-linear on the opposite strand. Taking the max of both
directions makes this metric strand-agnostic without needing to know the winning strand ahead of
time.

\section{\code{match\_rest}: best and second-best search}
\label{sec:match-rest}

\paragraph{Input.} All of a read's sorted, surviving buckets; the read's seed list, seed table,
and \code{diff\_hist}; an acceptance threshold \code{thr}; an optional ``forbidden'' \code{Mapping}
(when searching for the second-best, this is the best mapping already found, so that a
near-duplicate of it is excluded, \S\ref{sec:overlap-exclusion}); and \code{max\_overlap} (CLI
\code{-o}).

\paragraph{Output.} The single best \code{Mapping} across all surviving buckets that both beats
\code{thr} and (if \code{forbidden} is given) does not overlap it by more than
\code{max\_overlap}; or none, if no bucket qualifies.

\begin{algobox}{title={\code{match\_rest}(p\_sz, m, lmax, seeds, buckets, sorted\_buckets, diff\_hist, seed\_table, thr, forbidden, max\_overlap, metric, k)}}
\begin{verbatim}
best := None
for (loc, content) in sorted_buckets:              # most-matches-first, Section 8.6
    sh := 1.0
    if seed_heuristic_pass(buckets, seeds, m, loc, content, &sh, thr):   # Chapter 9
        candidate := findBestMapping(buckets, loc, content, seed_table, diff_hist,
                                      p_sz, m, lmax, sh, metric, k)
        if candidate.score() > thr:
            if forbidden is None or overlap(candidate, forbidden) < max_overlap:
                thr := candidate.score()            # raise the bar for subsequent buckets
                best := candidate
return best
\end{verbatim}
\end{algobox}

\subsection{Overlap exclusion}
\label{sec:overlap-exclusion}

\[
  \mathrm{overlap}(a, b) = \begin{cases} -0.0 & \text{different segments} \\
  \dfrac{\max(0,\ \min(a.t_r,b.t_r) - \max(a.t_l,b.t_l))}{\max(a.t_r,b.t_r) - \min(a.t_l,b.t_l)}
  & \text{same segment} \end{cases}
\]
i.e., intersection-over-union of the two candidates' matched reference intervals, restricted to
same-segment pairs. A second-best candidate search rejects any bucket whose refined interval
overlaps the already-found best by $\ge \mathit{max\_overlap}$ --- without this, the ``second
best'' would routinely just be a slightly-shifted restatement of the exact same true alignment
(since neighboring, heavily-overlapping buckets, per \S\ref{sec:bucket-geometry}, often all score
similarly well against the same true locus), which would make MAPQ (Chapter~\ref{ch:mapq})
uselessly report low confidence for essentially every uniquely-mapping read.

\subsection{How the caller invokes \code{match\_rest} twice}

The per-read driver (fully assembled in Chapter~\ref{ch:end-to-end}) calls \code{match\_rest}
exactly twice: once with \code{thr}$=\theta$ (the CLI acceptance threshold) and
\code{forbidden}$=$\code{None} to find the best mapping; then, only if a best mapping was found,
a second time with \code{thr} $= \mathit{best.score()} \times (1 - \mathit{min\_diff})$ and
\code{forbidden}$=$\code{best} to find the second-best. Both calls sweep the \emph{same} sorted
bucket list and share the \emph{same} \code{diff\_hist} (which is why every refinement sweep is
required to leave \code{diff\_hist} in exactly the state it found it, \S
\ref{sec:scoring-sliding-window} --- the second call depends on that invariant holding after the
first).

\section{The ground-truth-only sibling functions, and a genuine bug}
\label{sec:matcher-best-fixed-length}

A second, narrower set of functions (grouped under a type usually called \code{Matcher}) exists
solely to support the \emph{verbose ground-truth diagnostic} output (\S\ref{sec:ground-truth},
used only with \code{-v~2}, never in the main scoring path above). It re-implements
\code{collect\_matches} (turn a bucket into an explicit list of \code{Match}es, rather than
sweeping the reference sketch directly), a \code{bestFixedLength} that sweeps that explicit list
using \code{hit.r <= l.hit.r + p\_sz} as its window bound (notably \emph{not} respecting an
\code{lmax} parameter it is nonetheless passed --- see the dead-parameter note in
Chapter~\ref{ch:defects}), and a \code{bestIncludedJaccard} that finds the best-scoring window
whose length falls within an explicit $[\ell_{\min}, \ell_{\max}]$ range (used to check, precisely,
what score the ground-truth region itself would have achieved).

\begin{bugbox}
\label{sec:defect-prevr}
\code{bestIncludedJaccard}'s two-pointer sweep scores \emph{mid-loop}, immediately after growing
the window's right boundary by one, using the position \emph{one before} the just-processed
element (an iterator idiom equivalent to \code{prev(r)}) for both the matched span and the
reported right boundary --- while \code{intersection} at that exact point \emph{already} reflects
the just-processed (current, not ``previous'') element. This is an internal inconsistency: the
score's span excludes an element that the match count includes, which can and does produce
\code{intersection > span}, violating the metric's own well-formedness invariant
($\mathit{intersection} \le s\_sz$ must hold for a containment/Jaccard ratio to stay in $[0,1]$).
The original source itself carries a \texttt{// TODO: should it be prev(r) instead?} comment on
this exact line, i.e.\ the original author flagged this specific spot as unsettled. Confirmed by
constructing a small worked example (a $\sim$50bp read against a $\sim$200bp synthetic reference)
that reliably triggers the invariant violation. Because this path has \emph{zero} automated test
coverage in the reference codebase and only ever affects a diagnostic tag (never the main PAF
output), there is no existing oracle to defer to; the internally consistent fix --- score using
the \emph{current}, just-processed right-boundary element throughout, matching what
\code{intersection} already counts, exactly as \S\ref{sec:best-fixed-length}'s ``score after the
loop exits'' pattern already does correctly --- is the recommended behavior for a from-scratch
implementation of this diagnostic feature. See Chapter~\ref{ch:defects} for the full write-up.
\end{bugbox}

\section{Complexity}

Each refinement sweep (\code{bestFixedLength} or \code{lcs}) costs $O(|\mathrm{bucket}|)$: one
linear pass over the bucket's span of the reference sketch, with $O(1)$ amortized work per
position via the \code{diff\_hist} device (\S\ref{sec:scoring-sliding-window}) or, for LCS,
$O(|\mathrm{bucket}|\log|\mathrm{bucket}|)$ owing to the patience-sorting binary search. Summed
over all buckets that survive pruning (Chapter~\ref{ch:pruning}), and since pruning is designed
to make that survivor count small and roughly independent of total reference size, total
refinement cost per read is, in practice, far closer to $O(\ell)$ (one bucket's worth of work)
than to $O(n)$ (the whole reference).
